\documentclass[11pt,a4paper]{article}
\usepackage[utf8]{inputenc}
\usepackage[T1]{fontenc}
\usepackage[portuguese]{babel}
\usepackage{hyperref}
\usepackage{geometry}
\geometry{margin=25mm}
\usepackage{parskip}

\title{\textbf{Capítulo 1: Primeiros passos}}
\author{Curso de DevOps com Kubernetes}
\date{}

\begin{document}

\maketitle

\tableofcontents
\vspace{1em}

\section{Pré-requisitos}
Espera-se que os participantes tenham concluído o \href{https://courses.mooc.fi/org/uh-cs/courses/devops-with-docker}{DevOps com Docker} ou tenham experiência com Docker e Docker Compose. Além disso, é necessária experiência com desenvolvimento web, como Desenvolvimento Web Full Stack ou equivalente.

\section{Material do curso}
O material do curso deve ser lido parte por parte, do início ao fim. O material contém exercícios, que estão posicionados de forma que o conteúdo anterior forneça informações suficientes para a resolução de cada um. Faça os exercícios à medida que avança no material. Conforme você progride, precisará buscar cada vez mais informações na internet. Você também deve complementar seu conhecimento com a \href{https://kubernetes.io/docs/home/}{documentação oficial}!

\section{Uso de LLMs}
Grandes modelos de linguagem (LLMs), como \href{https://chatgpt.com/}{ChatGPT}, \href{https://claude.ai/new}{Claude} e \href{https://github.com/features/copilot?locale=pt-br}{GitHub Copilot}, provaram ser muito úteis no desenvolvimento de software. Pessoalmente, utilizo principalmente o Copilot, que se integra perfeitamente ao VS Code graças a um \href{https://visualstudio.microsoft.com/pt-br/github-copilot/}{plugin} e também se mostrou bastante útil no contexto do Docker.

Existem muitos casos de uso para os LLMs. Por exemplo, você pode:
\begin{itemize}
    \item Pedir ao LLM por várias configurações
    \item Pedir ao LLM para esclarecer o que uma determinada configuração faz exatamente
    \item Colar mensagens de erro e configurações no LLM e perguntar onde pode estar o problema
    \item Pedir esclarecimentos sobre conceitos fundamentais
\end{itemize}

O grau de utilidade das dicas fornecidas pelo Copilot e por outros modelos de linguagem varia. Talvez o maior problema desses modelos seja a \href{https://en.wikipedia.org/wiki/Hallucination_(artificial_intelligence)}{alucinação}: eles às vezes geram respostas que parecem muito convincentes, mas que estão completamente erradas. Ao programar, é claro que o código "alucinado" costuma ser detectado rapidamente quando não funciona. As situações mais problemáticas ocorrem quando o código ou a configuração gerada pelo modelo parece funcionar, mas contém bugs difíceis de detectar, como vulnerabilidades de segurança.

Outro problema na aplicação de modelos de linguagem ao desenvolvimento de software é a dificuldade que eles têm para "entender" projetos grandes e, por exemplo, gerar funcionalidades que exijam mudanças em vários arquivos. Atualmente, os modelos de linguagem também não conseguem generalizar o código; ou seja, se o seu código tiver funções ou componentes existentes que o modelo poderia usar com pequenas alterações, ele não fará isso. O resultado pode ser a deterioração da sua base de código, pois os modelos de linguagem geram muita repetição. Veja mais sobre isso, por exemplo, \href{https://visualstudiomagazine.com/articles/2024/01/25/copilot-research.aspx}{aqui}.

Ao usar modelos de linguagem, a responsabilidade é sempre do programador.

O rápido desenvolvimento dos modelos de linguagem coloca o aluno em uma posição desafiadora: vale a pena (ou sequer é necessário) aprender as coisas num nível detalhado quando se pode conseguir quase tudo pronto através dos modelos de linguagem? Nesse ponto, vale a pena lembrar da velha sabedoria de \href{https://en.wikipedia.org/wiki/Brian_Kernighan}{Brian Kernighan}, coautor do livro \textit{The C Programming Language}:

\begin{quote}
``Depurar é duas vezes mais difícil do que escrever o código em primeiro lugar. Portanto, se você escreve o código da forma mais inteligente possível, você, por definição, não é inteligente o suficiente para depurá-lo.'' --- Brian Kernighan
\end{quote}

Em outras palavras, como a depuração é duas vezes mais difícil do que a programação em si, não vale a pena programar um código que você mal consegue entender. Como seria possível depurar algo se a programação for terceirizada para um modelo de linguagem e o desenvolvedor de software não entender absolutamente nada do código? O mesmo se aplica à configuração do Kubernetes, e pode ser até mais grave: se você não entende os fundamentos, depurar problemas complexos do Kubernetes é simplesmente impossível.

Até agora, o desenvolvimento dos modelos de linguagem e da inteligência artificial continua numa fase em que eles não são autossuficientes, deixando os problemas mais difíceis para os humanos resolverem. Devido a isso, mesmo desenvolvedores iniciantes devem aprender os fundamentos muito bem, por precaução. Pode ser que, apesar da evolução dos modelos, conhecimentos ainda mais aprofundados se façam necessários. A inteligência artificial faz as tarefas fáceis, mas os humanos precisam resolver as confusões mais complicadas causadas pela IA. O GitHub Copilot tem um nome muito apropriado: ele é um Copiloto, um segundo piloto que auxilia o piloto principal em uma aeronave. O programador ainda é o piloto principal, o capitão, e carrega a responsabilidade final.

\section{Sobre o material}
O curso DevOps com Kubernetes foi criado pelo estimado Docker Captain \href{https://www.docker.com/contributors/jami-kousa/}{Jami Kousa} com a ajuda do departamento de ciência da computação da Universidade de Helsinque, pela academia de desenvolvimento de software (\href{https://toska.dev/}{Toska}). A partir de 2024, o curso passou a ser liderado por Matti Luukkainen, o criador do Full Stack Open. Você pode contatá-lo pelo e-mail \href{mailto:matti.luukkainen@helsinki.fi}{matti.luukkainen@helsinki.fi}.

Este material está licenciado sob a Licença \href{https://creativecommons.org/licenses/by-nc-sa/3.0/}{Creative Commons BY-NC-SA 3.0}, portanto você pode usá-lo e distribuí-lo livremente, desde que os criadores originais recebam os créditos. Se você fizer alterações no material e quiser distribuir a versão alterada, ela deverá ser licenciada sob a mesma licença. O uso deste material para fins comerciais é proibido sem permissão.

\end{document}